De Abdis verwoordde
waarom de beul moordde.
De gestrafte hoorde

bij de kwakkwajongens,
die voor hun genoegens
en soms wraakgevoelens

Emmy steeds belaagden.
Haar ter sluiks opjaagden
en het dan kloek waagden

te falen als naaiers.
De sneue zelfaaiers
verwerden tot zaaiers

van hun erfelijk spul
met enkel slap gelul.
Hij bleek geen lepe knul

want hij werd, vuig verklapt,
als enige betrapt.
Hij werd eerder gesnapt

veroordeeld tot schoolgang.
Daar hield een vrome tang
hem vast wel in bedwang.

Nou dat moest dan blijken.
Ze ging over lijken
en alles moest wijken

als de Abdis langskwam.
Messing was de stormram,
die de jongen aannam.
Ze zou hem opvoeden.
Hem voor kwaad behoeden.
Ze kon niet bevroeden

dat hij haar invloedssfeer
ontvluchtte keer op keer.
Ze schreef dan zijn naam weer

op het grote kalkbord.
Niet te zijn opgesnord
was voor hem juist een sport.

Het mijden van zijn lust
bracht alleen maar onrust.
Wat was, werd niet gesust.

Daardoor verloor Emmy
zich in de creatie
die uit de fantasie

kwam, hoewel veel vager,
van de mensenjager.
Haar geest werd voordrager

van haar eigen versie,
maar zag niet zijn fictie.
De kracht van illusie

hielp haar niet overeind.
Emmy bleef aangelijnd.
Zijn lust werd doorgeseind.

Zonder bewust te zijn
ontving Emmy’s brein
zijn primitieve sein.

Zijn vuur kaapte haar geest.
Ze was eerst scherp geweest.
Ook al was ze bevreesd,

ze leek niet te zwichten.
Er stroomden inzichten
die haar leed verlichtten

en haar ruimte boden.
Maria had noden,
maar had ze verboden.

Dat had Emmy verblind.
“Ik hoop dat ik rust vind
als ik mijn lijf verbindt

aan deze oude man.
Dat is wat ik doen kan!
Wat moet ik anders dan?”

Ze miste reflectie
en werd haar perceptie.
Daarmee verdween Emmy.

Een milde lerares
sprong voor hem op de bres
toen ooit tijdens haar les

zijn vader binnenkwam.
Hij was in vuur en vlam.
Omdat de non aannam

dat zijn geprikkeldheid
over afwezigheid
zou leiden tot een strijd,

zei ze jachtig en kort:
“Hij staat niet op het bord!”
Er werd geenszins gemord,

geschreeuwd of gezanikt.
maar zelfs verheugd geknikt.
Het werd zomaar geslikt

en zelfs thuis gehuldigd.
Nu eens niet beschuldigd,
werd er taart genuttigd.

Men hoorde luid gevloek
bij Messings werkbezoek.
De jongen was weer zoek

en stond opgeschreven.
Zijn schuld opgeheven
duurde dus maar even.

Want wat bleek er later?
Men maakte een flater.
Het bord was met water

door de jongen geschoond
om te worden verschoond.
Wel was hij al beloond.

Er werd al vier weken
naar hem uitgekeken
tot zij hem kon spreken.

“Het spijt me zeer, Abdis.
Weet u wat het plat is.
Ik had het gewoon mis.

Ik meende dat u zei:
je hebt vier weken vrij.
Niet: de fier kweken blij,

toen u Spreuken besprak.”
De Abdis had echt lak
aan zijn valse aanpak.

Tot hij plots citeerde
wat hij van haar leerde:
“Jezus formuleerde

ieder te vergeven
die onwetend leven.”
Ze liet hem maar even.

Eens tijdens een pestbui
gaven de jongelui
aan hun schaamte de brui.

Eentje trok ineens vlug
al liggend op zijn rug
zijn onderbroek terug

als richtinggevende
toen ze haar ketenden.
Anderen openden

de rits van hun werkbroek.
Het ongevraagd bezoek
hield alvast een zakdoek

klaar voor de eruptie.
Het bekoorde Emmy
dat haar expositie

zou gaan leiden tot pulp,
die de klare zelfhulp
zichtbaar in elke gulp

uit de harde staken
zou gaan veroorzaken.
Ze mocht het meemaken.

De spanning opwekken.
De spanning aftrekken.
De spanning wegtrekken.

Al stal men wel een kus,
geen streelde als bonus
een deel van haar corpus.

Zo’n onkiese streling
zou ze bij vergeving
voor deze beleving

van moeten kokhalzen.
Nu kon ze reikhalzen
naar deze jakhalzen.

“Bepotelen die staak,
die ik langzaam ontwaak,
terwijl ik het aanraak.

Ik wil het vastgrijpen.
Verklaren hoe knijpen
hem stilaan laat rijpen.

Openbaren welk vuur
ontbrandt na korte duur
uit dit gestaafd figuur.

Vergaren het witgoed
dat blijkbaar zo goed doet
als het zijn eind ontmoet.”

Deze revelatie
over haar sensatie
smoorde haar reflectie.

Plompverloren greep hij
Emmy ferm bij haar zij.
Ze kwam met de schrik vrij.

Met slechts zijn linkerhand
draaide hij haar galant
uit de inspectiestand

en dwong haar te zitten.
Emmy wilde spitten,
om hem te verhitten,

in zijn doelstellingen.
Om wat voorspellingen
over haar kwellingen

zo te achterhalen.
Werden het schandpalen
of moest ze betalen

met haar jonge leven?
De kou liet haar beven
net als dat gegeven.

Haar ontblote billen
begonnen te trillen.
“Wat zou de man willen?”,

ging steeds door Emmy heen.
Hij keek kalm om zich heen
met een hand op haar been,

beschouwend het decor.
Zijn hut als monitor.
De ven voor haar comfort.

De ven waar zoeven
een bovenaards wezen
voor hem was verrezen.

“Met het ingetogen
mystieke vermogen
het licht van mijn ogen

innig te beminnen
via mijn rijmzinnen.
Hoe kan ik beginnen

mijn geest op te beuren,
door haar te besmeuren
zonder te verkleuren

haar beeldige schoonheid?
Hoe houdt ze haar echtheid
en haar ware blijheid

als ik haar corrigeer
en haar manipuleer?
Hoe blijft haar goede eer,

als ik haar belazer?
Krijgt ze op haar lazer,
wordt ze dan niet dwazer?

Wordt ze afgepeigerd,
verhandeld, vernederd
en besodemieterd,

raakt ze dan voor altijd
haar eerlijke grootheid
en eenheid aan mij kwijt.

Kan ik haar opvoeden
als wel beïnvloeden
en tevens behoeden

haar zuivere kindsheid,
haar montere speelsheid
en haar lichte vrijheid.”

Ik gaf bedwelmd uiting
aan deze mijmering
middels zijn fluistering.

Waardoor Emmy die zweeg
het volledig meekreeg.
Haar ongerustheid steeg.

Ze moest zachtjes zuchten.
Wat had ze te duchten?
Kon ze beter vluchten?

Hij las haar gedachten
en wist uit zijn jachten
wat hij kon verwachten.

Hij maakte eerder mee
dat plots een jachttrofee
zich als prostituee

wist te presenteren.
Om hem dan te smeren,
zonder uit te keren

de prijs voor de vrijheid
bij onoplettendheid.
“Niet bij deze schoonheid”,

broedde de drinkebroer.
Hij griste de halssnoer
van de aspirant hoer

haar nog vochtige nek.
Haar vreemde moedervlek
was de finale check

door de romanticus.
“Met Emmy Nativus
heb ik te maken dus!

Je wordt schijnbaar gezocht.”
Emmy verloor traanvocht.
Ze zal worden verkocht

als een gevangen prooi.
Te leven in een kooi
voor iemands geflikflooi

of misschien zelfs de dood
verschafte hoge nood.
Of Emmy zich aanbood?
